\documentclass[10pt,a4paper]{article}
\usepackage{amsmath,amssymb,amsthm,graphicx}
\usepackage{hyperref}
\usepackage{cleveref}
\usepackage{amssymb,amsmath,amsthm}
\theoremstyle{plain}
\newtheorem{theorem}{Theorem}
\newtheorem{lemma}[theorem]{Lemma}
\newtheorem{observation}[theorem]{Observation}
\theoremstyle{definition}
\newtheorem{remark}[theorem]{Remark}
\title{A large subset meeting every large edge in zero or two vertices, found in cubic time}
\author{Johan Hastad, Bjorn Martinsson, Tamio-Vesa Nakajima, Stanislav Zivny}
\date{arxiv:2404.19556v6; Creative Commons Attribution 4.0 International}
\begin{document}
\maketitle
\noindent\textit{Editorial note.} The counted statement is the author's labelled environment at source0.tex lines 151-156 and its proof environment opens at line 158; the pairing is by adjacency and the retained proof range is the whole environment.\par\medskip
\noindent\textit{Editorial note.} This source is newly approved for this line and its preamble is built from the author's own package line and theorem-environment family (source lines 5-23); the font, language and layout packages of the source are not retained because the retained material does not use them.\par\medskip
\noindent\textit{Editorial note.} Metadata provenance: the pinned source carries a licence evidence file recording source\_version 2404.19556v6, CC BY 4.0, and a separate source record file; the card's source\_version is the value recorded in the licence evidence.\par\medskip
\noindent\textit{Editorial note.} The source carries no label-number dump and no counter values are imposed, so the printed numbering is the wrapper's default; identity is the author label.\par\medskip
\noindent\textit{Editorial note.} r002 note: r001 was refused by the assembler because the card carried an empty author list; r002 fills in the authors and the source title recorded in the source record of the pinned copy and changes nothing else.\par\medskip
\noindent\textit{Editorial note.} r003 note: the pre-publication validator requires either status=supported in the licence evidence or a retained-licence form with no status at all. This source's evidence file carries no acceptance record and none may be written on anyone's behalf, so r003 keeps no status and records a retained source notice (source\_notice), quoting the author's own CC BY declaration on lines 47-48 of the manuscript with its path form in a task-directory evidence copy that the assembler carries into the package: the manuscript declares CC BY at lines 47-48. The original keys of the source evidence are retained unchanged, no new licence claim is made, and no acceptance record is written.\par\medskip
\noindent\textit{Editorial scope.} One original result of the pinned author manuscript, counted once: the inner step producing a large subset meeting every triple in zero or two vertices (author label lem:innerstep).\par\medskip
\noindent\textit{Allowed standard prerequisites (editorial).}\begin{itemize}
\item Linearly ordered colouring. A colouring of the vertices such that in every edge the colour that appears once and is smallest among the colours present is unique.. It is the hypothesis on the hypergraph.
\item Implied linear system. The colouring conditions translate into a system of linear equations modulo two over the vertex set, and no variable is fixed when the solution space is not contained in a coordinate hyperplane.. It is the algebraic description the algorithm works with.
\item Basic solution and parametrisation. The solution set is an affine space, so a generic solution is a basic solution plus a combination of basis vectors.. It is what makes the randomised choice of coefficients possible.
\end{itemize}\medskip
\begin{lemma}\label{lem:innerstep}
    There is an algorithm which, if given a $\{2, 3\}$-uniform hypergraph $H$ with $n$ vertices and
    $m$ edges that admits an LO 2-colouring and such that the implied linear
    system of equations $Av=1^m$ does not fix the value of any variable, outputs a subset $T$ of vertices that intersects edges of size three in zero or two vertices and edges of size two in exactly one vertex.
    Moreover, we have  $\vert T\vert \geq n / 2$. The algorithm runs in $O(n^3 + nm)$ time.
\end{lemma}

\begin{proof}

    We first describe a randomised version of our algorithm, and then derandomise
    it. The set of solutions to $Av=1^m$ is an affine space and hence a generic solution can be written as $v=v^0+\sum_{i=1}^r a_i v^i$ for a basic solution $v^0$, linearly independent solutions to the homogeneous system $v^i$, and field elements (in this case bits) $a_i$.
    The fact that no variable is fixed implies that for each vertex $x$ there is 
    some positive $i$ such that $v^i_x=1$.
    
     For the randomised algorithm choose $a_1, \ldots, a_r$ to be independent identically distributed uniformly random bits, and set $T$ to be 
     the set of variables $x$, such that $v_x=0$.  Clearly $T$ satisfies
     the conclusion of the lemma as in each edge we have an odd number of ones. 
     For every vertex $x$, there exists positive $i$ such that $v_x^i = 1$ --- hence, due to the influence of $a_i v_x^i$, we see that $x$ is included in $T$ with probability $\frac{1}{2}$.
     Thus, on average $T$ contains half the vertices.

    Now, we derandomise this algorithm using the method of conditional
    expectations. Go through the variables $a_i$ in increasing order and 
    fix its value once and for all.  Fixing the value of $a_i$ determines
    the value(s) of some $v_x$ while other values remain undetermined.  For each
    value being determined $v_x^i=1$ and hence one value of $a_i$ gives
    the final value 0 and the other gives final value 1.  Set $a_i$ such
    that at least half the determined values are 0.  After we have 
    fixed all $a_i$ this way, we have a final solution with at least
    $n/2$ zeroes.
    
    The bottleneck of the running time of this algorithm is solving the linear system of equations. This can be done in the advertised running time since every equation has $O(1)$ entries. 
\end{proof}

\end{document}
